% Generated draft: requires mathematical and source review.
\documentclass[11pt]{article}
\usepackage[margin=25mm]{geometry}
\usepackage{fontspec}
\setmainfont{Latin Modern Roman}
\usepackage{amsmath,amssymb,mathtools,unicode-math}
\setmathfont{Latin Modern Math}
\usepackage{enumitem,needspace,xurl,hyperref}
\setlength{\parindent}{0pt}
\setlength{\parskip}{5pt}
\setlength{\emergencystretch}{3em}
\newcommand{\N}{\mathbb{N}}
\newcommand{\Z}{\mathbb{Z}}
\newcommand{\Q}{\mathbb{Q}}
\newcommand{\R}{\mathbb{R}}
\newcommand{\C}{\mathbb{C}}
\newcommand{\F}{\mathbb{F}}
\newcommand{\A}{\mathbb{A}}
\newcommand{\PP}{\mathbb{P}}
\newcommand{\E}{\mathbb{E}}
\newcommand{\eps}{\varepsilon}
\newcommand{\dd}{\,\mathrm{d}}
\begin{document}
\section*{UnsolvedMath ID: 30001039}
\textbf{Source problem code:} catalogue entry
\subsection{Definitions and mathematical statement}
Let \(G\) be a Lie group with real Lie algebra \(\mathfrak g\), and let \(K\subseteq\mathfrak g\) be the adjoint-invariant cone in the source. A Lie quasi-state in the stated algebraic sense is a function \(\zeta:\mathfrak g\to\R\) such that
\[
 [X,Y]=0\quad\Longrightarrow\quad
 \zeta(aX+bY)=a\zeta(X)+b\zeta(Y)
 \quad(a,b\in\R),
\]
with \(\zeta(X)\geq0\) for \(X\in K\). Equivalently, its restriction to every commutative linear subspace is linear. The general request is to detect or characterize all such functions for the given \((\mathfrak g,K)\). Adjoint invariance of \(K\) is a hypothesis on the cone; it does not itself impose \(\zeta(\operatorname{Ad}(g)X)=\zeta(X)\). Continuity, measurability, local boundedness, and invariance of the quasi-state are not stated.

The source then asks three groups of special questions. First, is the Maslov Lie quasi-state for \(\operatorname{Sp}(2n)\) unique up to the intended isomorphism? The usual setting is the real symplectic group and its Lie algebra, but the cone, normalization, range of \(n\), regularity class, and meaning of isomorphism are absent. Positive scaling and pullback by automorphisms are different equivalences; none is selected.

Second, does \((T^4,\omega_{\mathrm{std}})\) admit a nonlinear Lie quasi-state? Third, how many do \((\mathbb{CP}^2,\omega_{\mathrm{std}})\) and \((S^2\times S^2,\omega_{\mathrm{std}}\oplus\omega_{\mathrm{std}})\) admit? A manifold itself is not a Lie algebra. The standard symplectic interpretation uses \(C^\infty(M,\R)\) with Poisson bracket and the cone of pointwise nonnegative functions, giving
\[
 \{f,g\}=0\quad\Longrightarrow\quad
 \zeta(af+bg)=a\zeta(f)+b\zeta(g).
\]
Nonlinear means failing to be linear on all functions, rather than on just the Poisson-commuting subspaces. The source does not specify whether this full Poisson algebra, functions modulo constants, or another Hamiltonian Lie algebra is intended, nor whether normalization \(\zeta(1)=1\), monotonicity, continuity, or symplectic invariance is additionally required. Definitions of symplectic quasi-states used in the cited work can include more axioms and must not be equated automatically with the minimal Lie quasi-state definition above.

The supplied assessment dated 2026-08-21 reports uniqueness and multiplicity results in important cases, including rank-at-least-three cases and \(S^2\times S^2\), but says full detection/classification and the \(T^4\) question remain unresolved. It supplies neither the regularity hypotheses nor exact counts or equivalence classes of those results.
\subsection{Short English statement}
Classify functions linear on commuting Lie-algebra elements and positive on a given invariant cone, then settle the stated Maslov uniqueness and symplectic-manifold existence/counting questions. Regularity, normalization, equivalence, and the precise Poisson-algebra interpretation are essential missing choices.
\subsection{Sources}
\begin{itemize}
\item Local UnsolvedMath registry, permanent ID \(30001039\), exact source code \texttt{catalogue entry}, proposed year 2008; no individual proposer is supplied.
\item \emph{Geometric Group Theory, Hyperbolic Dynamics and Symplectic Geometry (2009)}, Oberwolfach Reports, DOI 10.4171/owr/2008/42, pp. 2337--2406. \url{https://doi.org/10.4171/owr/2008/42}.
\item Entov--Polterovich, ``Lie quasi-states,'' \url{https://arxiv.org/abs/0906.4901}.
\item Eliashberg--Polterovich, ``symplectic quasi-states,'' \url{https://arxiv.org/abs/1006.2501}.
\item The supplied \texttt{corrected\_verified} note reports preservation of the full question with notation cleanup. The partially solved label and assessment checked 2026-08-21 are supplied dated curation, not independently verified; exact rank, regularity, normalization, and manifold hypotheses require review.
\end{itemize}
\end{document}
